\section{Writing Conventions} 
In general I use British English.
\subsection{Logical rules}
Each item in a list will end with a point and start capitalized.
If I list something where the order matters like:
\begin{enumerate}
	\item Step one.
	\item Step two.
	\item Step three.
\end{enumerate}
and for equivalent properties I use letters: For example, a Morse-Smale function $f$ on a Riemannian manifold $(M,g)$ satisfies:
\begin{enumerate}[(a)]
	\item $f$ is smooth.
	\item The critical points of $f$ are non degenerate.
	\item The (un-) stable manifolds intersect transversal.
\end{enumerate}

\subsection{Structure}
In this thesis i have the following text-types
\begin{enumerate}[(a)]
	\item Free text that is not in any container.
	\item The theorem container.
	\item The lemma container.
	\item The corollary container.
	\item The proposition container.
	\item The prove container.
\end{enumerate}
Free Text is the main narrative of the thesis. It includes explanations, motivations, intuitions, examples, informal remarks, and transitions. Free text connects the formal content and helps guide the reader through the logical structure of the work.
\begin{theorem}[An Example]
	Used for major, central mathematical statements that are either original contributions or critical results from the literature. Theorems typically capture the most important claims and are often proved in detail.
\end{theorem}
\begin{lemma}[An Example]
	Used for technical results that support the proof of theorems or other key statements. Lemmas are usually not of independent interest but are essential building blocks within a logical argument.
\end{lemma}
\begin{corollary}[An Example]
	Used for statements that follow immediately and straightforwardly from a previously proven theorem or lemma. Corollaries are usually included to highlight direct consequences of major results.
\end{corollary}
\begin{proposition}[An Example]
	Used for results that are mathematically relevant but not as central as theorems. Propositions often encapsulate useful facts, structural properties, or classifications that are too substantial for a lemma but not significant enough for a theorem.
\end{proposition}
\begin{proof}
	This container is pretty obvious
\end{proof}
\subsection{Mathemtical Conventions}
Elements in a vector bundle will have greek letters. For example $\xi\in TM$	